\documentclass[11pt,a4paper]{article}

\usepackage[margin=1.65cm]{geometry}
\usepackage{titlesec}
\usepackage{enumitem}
\usepackage{hyperref}
\usepackage{xcolor}

% Colors
\definecolor{linkcolor}{RGB}{0,70,130}

% Hyperref setup
\hypersetup{
    colorlinks=true,
    linkcolor=linkcolor,
    urlcolor=linkcolor
}

% Section formatting
\titleformat{\section}{\large\bfseries\uppercase}{}{0em}{}[\titlerule]
\titlespacing{\section}{0pt}{6pt}{2pt}

% Remove page numbers
\pagenumbering{gobble}

% Slightly tighter leading so the longer experience section stays on one page
\linespread{0.95}

% Custom commands
\newcommand{\entry}[4]{
    \textbf{#1} \hfill #2 \\
    \textit{#3} \hfill \textit{#4}
}

\begin{document}

% Header
\begin{center}
    {\LARGE\bfseries Vinícius Bregoli} \\[4pt]
    KI/ML-Ingenieur \\[1pt]
    Computer Vision \enspace$\cdot$\enspace LLM-Infrastruktur \enspace$\cdot$\enspace KI-Agenten \\[6pt]
    \href{https://www.viniciusbregoli.dev}{viniciusbregoli.dev} \enspace$\cdot$\enspace
    \href{mailto:me@viniciusbregoli.dev}{me@viniciusbregoli.dev} \enspace$\cdot$\enspace
    Curitiba, Paraná \\[2pt]
    \href{https://linkedin.com/in/viniciusbregoli/}{linkedin.com/in/viniciusbregoli} \enspace$\cdot$\enspace
    \href{https://github.com/viniciusbregoli}{github.com/viniciusbregoli}
\end{center}

\section{Technische Kenntnisse}
\noindent\textbf{Programmiersprachen:} Python, Rust, JavaScript, C/C++, SQL, HTML, CSS \\
\textbf{Frameworks und Bibliotheken:} FastAPI, Node.js, PyTorch, OpenCV, YOLO \\
\textbf{Tools und Plattformen:} Git, Linux, PostgreSQL, Redis, Virtualisierung \\
\textbf{Bereiche:} Computer Vision, 3D-Inferenz, Machine Learning, LLM-Infrastruktur, KI-Agenten (MCP), Full-Stack-Entwicklung, Deployment, IT-Infrastruktur

\section{Berufserfahrung}

\entry{ML/KI-Entwickler}{Kinebot}{Curitiba, Brasilien}{Okt 2025 -- Heute}
\begin{itemize}[nosep, leftmargin=*]
    \item Entwicklung von Computer-Vision-Modellen (YOLO) und 3D-Inferenz zur Unterstützung ergonomischer Videoanalysen.
    \item Aufbau eines Agenten-Orchestrators für die ergonomische Videoanalyse: spezialisierte Agenten (\mbox{Vision} sowie die Bewertungen REBA, KIM-PP und NIOSH) laufen parallel über einen vom Orchestrator verwalteten gemeinsamen Kontext.
    \item Entwicklung eigener MCP-Server für die internen Tools des Unternehmens.
    \item Integration generativer KI-Modelle in die ergonomischen Analysen der Plattform.
    \item Zeitanalyse von Videos: Identifikation von Aktivitäten und Einordnung ihrer Start- und Endintervalle in eine Zeitleiste.
    \item GPU-Infrastruktur für das Serving von LLMs.
    \item Python-Backend-Dienste (FastAPI, Redis, PostgreSQL), die die KI-Verarbeitung mit den Workflows der Plattform verbinden.
\end{itemize}

\vspace{3pt}

\entry{Praktikant -- Softwareentwicklung}{Exati Tecnologias}{Curitiba, Brasilien}{2022 -- 2023}
\begin{itemize}[nosep, leftmargin=*]
    \item Entwicklung und Wartung von Webanwendungen mit Vue.js
    \item Integration von Frontend-Anwendungen mit Backend-Diensten mittels Java und Node.js
    \item Entwurf und Verwaltung von PostgreSQL-Datenbanken
    \item Zusammenarbeit mit Git/GitLab zur Versionskontrolle und YouTrack zur Aufgabenverwaltung
\end{itemize}

\vspace{3pt}

\entry{Praktikant -- IT-Infrastruktur}{Siemens Energy Brasil}{Curitiba, Brasilien}{2021 -- 2022}
\begin{itemize}[nosep, leftmargin=*]
    \item Verwaltung von IT-Infrastruktur und Virtualisierungsumgebungen
    \item Remote-Support in einem anspruchsvollen Unternehmensumfeld
    \item Bearbeitung von Anfragen und Ticket-Management
\end{itemize}

\section{Ausbildung}

\entry{Informatik und KI (Austauschprogramm)}{Technische Hochschule Ingolstadt}{Ingolstadt, Deutschland}{Okt 2024 -- Aug 2025}

\vspace{3pt}

\entry{B.Eng. Technische Informatik}{Pontifícia Universidade Católica do Paraná (PUCPR)}{Curitiba, Brasilien}{2020 -- 2025}

\section{Projekte und Forschung}

\textbf{Digital-Twins-Forschung -- Airbus-Partnerschaft} \hfill 2021 \\
Stipendienprogramm für technologische Entwicklung und Innovation mit Schwerpunkt auf Digital Twins und deren Anwendungen in der Luft- und Raumfahrt.

\vspace{3pt}

\textbf{PUCPR Mobile Robotics Team} \hfill 2020 -- 2023 \\
Mitwirkung bei der Entwicklung von Linienfolger-, Sumo-, Trekking- und Kampfrobotern. Teilnahme an Wettbewerben wie RoboCore Experience (RCX) und IRONCup.

\section{Sprachen}

\noindent Englisch (Fließend) \enspace$\cdot$\enspace Portugiesisch (Muttersprache) \enspace$\cdot$\enspace Spanisch (Mittelstufe) \enspace$\cdot$\enspace Deutsch (Grundkenntnisse)

\end{document}
